\section{Por qué la detección de vulnerabilidades es difícil: contexto semántico, ruido en la evaluación y verificación dinámica}

\subsection{Caso del kernel: lo sospechoso a nivel local no implica que sea válido a nivel global}
El curso usó un ejemplo de código del kernel para explicar la dificultad de determinar vulnerabilidades: una función lee una longitud del encabezado de un paquete, asigna un búfer y copia datos; si la longitud no coincide con el payload real, puede producirse una escritura fuera de límites. La dificultad no está en el reconocimiento de patrones en sí, sino en \textbf{confirmar la ruta alcanzable, las restricciones del llamador y el flujo de datos entre funciones}.

\begin{figure}[H]
\centering
\includegraphics[width=0.88\textwidth]{frame-05-kernel-vuln.jpg}
\caption{Fragmento de la explicación de la vulnerabilidad del kernel: el riesgo de la misma función depende del contexto global de la llamada \protect\footnotemark}
\end{figure}
\footnotetext{Intervalo de tiempo en el video: 01:06:30--01:07:10.}

\begin{importantbox}{Juicio clave del ponente}
\textit{``This function on its own is not vulnerable or non-vulnerable.''}  
Clasificar una función de forma aislada suele no ser un problema bien definido. Que una vulnerabilidad se considere válida depende de si el llamador ya restringe la entrada, de si la ruta es alcanzable y de si el entorno de ejecución cumple las condiciones para desencadenarla.
\end{importantbox}

\subsection{La dificultad natural de evaluar con bases de datos de vulnerabilidades}
Las bases de datos públicas de vulnerabilidades (como el ecosistema CVE) parecen una fuente de datos ideal: tienen número, descripción y parche. Pero convertirlas directamente en una ``tarea de clasificación a nivel de función'' genera problemas serios: hasta dónde cortar el contexto de entrada, cómo definir los ejemplos positivos y negativos, si el parche corresponde a una sola causa raíz, y si la etiqueta es estable entre versiones.

\begin{warningbox}{``Descargar una base de datos de vulnerabilidades para entrenar un detector'' suele ser demasiado optimista}
La clase señaló explícitamente que esta es la ruta más natural, pero también la más fácil de errar. Si no se define primero un límite riguroso de la tarea, el modelo puede aprender el estilo del proyecto, los hábitos de los commits o frases de plantilla, en lugar de la semántica real de la explotabilidad.
\end{warningbox}

\subsection{La necesidad de una evaluación dinámica}
Sutton mencionó las competencias relacionadas con DARPA y la dirección de evaluación de seguridad de Meta; el punto en común es pasar de etiquetas estáticas a verificación ejecutable. Para la tarea de descubrimiento de vulnerabilidades, la evaluación dinámica se acerca más al objetivo real: activación por entrada, evidencia de ejecución, ubicación del fallo, reproducibilidad.

\begin{knowledgebox}{La relación complementaria entre fuzzing y los agentes}
\begin{itemize}
    \item El fuzzing es bueno para la exploración aleatoria de alto rendimiento;
    \item El agente es bueno para guiar por semántica y plantear hipótesis de ruta;
    \item Combinar ambos puede formar una búsqueda híbrida de ``cobertura amplia + ruta profunda''.
\end{itemize}
Esta es también la dirección que se enfatiza repetidamente en la última parte de la clase: la IA no sustituye al fuzzing, sino que empuja la búsqueda hacia las ramas lógicas profundas a las que el fuzzer difícilmente llega.
\end{knowledgebox}

\subsection{De la precisión estática a la tasa de éxito de la tarea}
En escenarios de seguridad, el indicador final no debería ser ``qué tanto se parece este código a una vulnerabilidad'', sino ``si se puede construir una entrada que la active y aportar evidencia auditable''. Esto traslada el foco de investigación de la clasificación de texto a \textbf{la tasa de éxito de tareas ejecutables}, algo que coincide de forma natural con el paradigma de agentes.

\begin{importantbox}{Implicación directa para el diseño de investigación}
Si en el futuro los artículos sobre agentes de seguridad solo reportan puntuaciones de clasificación estática, será difícil sustentar su valor de ingeniería; como mínimo deberían añadir indicadores como la tasa de activación dinámica, la estabilidad de reproducción y el costo de tratar los falsos positivos.
\end{importantbox}

\subsection{Resumen de la sección}
La dificultad de la detección de vulnerabilidades está en ``la dependencia del contexto más la verificación ejecutable'', no en ``la falta de reconocimiento de patrones''. Esto también explica por qué la seguridad necesita más sistemas basados en agentes, y no solo modelos de texto de código más potentes.
